\subsection{Métodos}

\subsubsection{Regla: PascalCase y verbo descriptivo}

\textbf{Regla:} El nombre de un método se escribe en \texttt{PascalCase} y comienza con un verbo o frase verbal (Microsoft, 2025c). Se exceptúan las convenciones definidas expresamente en este estándar y los nombres impuestos por una interfaz o por la plataforma .NET.

\textbf{Código correcto:}
\begin{lstlisting}[style=csharpstyle]
public int CalculateTotalScore(int basePoints, int bonusPoints)
{
    return basePoints + bonusPoints;
}
\end{lstlisting}

\textbf{Código incorrecto:}
\begin{lstlisting}[style=csharpstyle]
public int totalScore(int basePoints, int bonusPoints)
{
    return basePoints + bonusPoints;
}
\end{lstlisting}

\subsubsection{Regla: máximo de parámetros}

\textbf{Regla:} Un método no recibe más de tres parámetros; con más datos de entrada, se encapsulan en un DTO.

\textbf{Descripción:} McConnell señala que una lista larga de parámetros incrementa la probabilidad de invocar el método con el orden o el valor incorrecto (McConnell, 2004).

\textbf{Código correcto:}
\textit{Archivo \texttt{MoveRequest.cs}:}
\begin{lstlisting}[style=csharpstyle]
public readonly record struct MoveRequest(Player Player, Tile Origin, Tile Destination);
\end{lstlisting}

\textit{Archivo \texttt{MoveExecutor.cs}:}
\begin{lstlisting}[style=csharpstyle]
public sealed class MoveExecutor
{
    public MoveResult ExecuteMove(MoveRequest moveRequest)
    {
        return new MoveResult(moveRequest.Player, moveRequest.Destination);
    }
}
\end{lstlisting}

\textbf{Código incorrecto:}
\begin{lstlisting}[style=csharpstyle]
public sealed class MoveExecutor
{
    public MoveResult ExecuteMove(Player player, Tile origin, Tile destination, int diceRoll, bool isDoubleTurn)
    {
        return new MoveResult(player, destination);
    }
}
\end{lstlisting}

\subsubsection{Regla: métodos booleanos como pregunta}

\textbf{Regla:} Un método que devuelve \texttt{bool} y expresa una pregunta se nombra con prefijo interrogativo (\texttt{Is}, \texttt{Has}, \texttt{Can}) (Martin, 2008). Se exceptúan patrones convencionales de .NET como \texttt{Try...} y miembros definidos por la plataforma o por una interfaz, como \texttt{Equals}.

\textbf{Código correcto:}
\begin{lstlisting}[style=csharpstyle]
public bool IsMoveValid(Tile origin, Tile destination)
{
    return destination.IsWithinBoard && !destination.IsOccupied;
}
\end{lstlisting}

\textbf{Código incorrecto:}
\begin{lstlisting}[style=csharpstyle]
public bool CheckMove(Tile origin, Tile destination)
{
    return destination.IsWithinBoard && !destination.IsOccupied;
}
\end{lstlisting}
